% Generated draft: requires mathematical and source review.
\documentclass[11pt]{article}
\usepackage[margin=25mm]{geometry}
\usepackage{fontspec}
\setmainfont{Latin Modern Roman}
\usepackage{amsmath,amssymb,mathtools,unicode-math}
\setmathfont{Latin Modern Math}
\usepackage{enumitem,needspace,xurl,hyperref}
\setlength{\parindent}{0pt}
\setlength{\parskip}{5pt}
\setlength{\emergencystretch}{3em}
\newcommand{\N}{\mathbb{N}}
\newcommand{\Z}{\mathbb{Z}}
\newcommand{\Q}{\mathbb{Q}}
\newcommand{\R}{\mathbb{R}}
\newcommand{\C}{\mathbb{C}}
\newcommand{\F}{\mathbb{F}}
\newcommand{\A}{\mathbb{A}}
\newcommand{\PP}{\mathbb{P}}
\newcommand{\E}{\mathbb{E}}
\newcommand{\eps}{\varepsilon}
\newcommand{\dd}{\,\mathrm{d}}
\begin{document}
\section*{UnsolvedMath ID: 11000252}
\textbf{Source problem code:} catalogue entry
\subsection{Definitions and mathematical statement}
Conjecture 4.4 states that every closed hyperbolic $4$-manifold $M$ has all Seiberg--Witten invariants zero. In the usual oriented smooth setting, for $b_2^+(M)>1$ this means
\[
 \operatorname{SW}_M(\mathfrak s)=0
 \quad\text{for every }\operatorname{Spin}^c\text{ structure }\mathfrak s,
\]
with the relevant invariant/insertions and homology orientation conventions fixed. Here $b_2^+$ is the dimension of a positive subspace for the intersection form. The excerpt does not state orientability, chamber choices for $b_2^+=1$, or a convention for $b_2^+=0$. Consequently the word all cannot be turned into a chamber-independent numerical statement in those cases without further source data.

The source's Proposition 4.5 says that a positive answer would imply that no closed hyperbolic $4$-manifold is symplectic. It cites Taubes [37] for nonvanishing on compact oriented symplectic $4$-manifolds with $b_2^+>1$. For the remaining symplectic case $b_2^+=1$, it uses hyperbolic signature zero to obtain $b_2(M)=2$, and residual finiteness of $\pi_1(M)$ to pass to a finite cover $p:M_1\to M$ of degree $d>1$. The pulled-back symplectic form is symplectic. Its outlined contradiction, if $b_2(M_1)=2$, uses positive Euler characteristic and
\[
 \chi(M_1)=d\chi(M),\qquad
 b_1(M_1)=2-d(2-b_1(M)).
\]
The displayed arithmetic requires an additional justification, such as injectivity of pullback on rational first cohomology, or sufficiently large covering degree, to obtain the claimed contradiction in every case. The source refers to an earlier volume formula not included here. It also notes that $b_2=0$ cannot be symplectic.

For contrast, the source recalls Thurston's symplectic construction for surface bundles over surfaces with both genera at least $2$, in the usual closed oriented bundle setting. It then announces a generalization of omitted Question 4.1 analogous to virtual fibering in dimension three, but gives no further question in this record. None of this context supplies a proof of the vanishing conjecture.
\subsection{Short English statement}
The source conjectures vanishing of all Seiberg--Witten invariants for closed hyperbolic $4$-manifolds and explains why that would obstruct symplectic structures. Orientation and low-$b_2^+$ chamber conventions need clarification.
\subsection{Sources}
\begin{itemize}
\item Local registry: \url{data/unsolvedmath/problems.json}, permanent UnsolvedMath ID $11000252$, exact source code \texttt{catalogue entry}.
\item Farb, \emph{Problems on Mapping Class Groups and Related Topics}, Conjecture 4.4 and contextual Proposition 4.5, PDF page 288, extracted footer A.W. Reid: \url{https://www.math.uchicago.edu/~farb/papers/mcgbook.pdf}. Supplied access date: 2026-07-29; extraction: pdf-font-structure; page-disambiguated numbering among 329 headings.
\item Supplied references: Taubes [37]; Thurston's argument via [39] or [14], Theorem 10.2.17. Full entries, the earlier volume formula, Question 4.1 and the following generalization are absent. The registry literature summary is truncated.
\end{itemize}
\end{document}
